\documentclass[12pt,a4paper]{article}
\usepackage[margin=1in]{geometry}
\usepackage{booktabs}
\usepackage{longtable}
\usepackage{graphicx}
\usepackage{hyperref}
\title{\textbf{Vero Cell Culture Morphology Detection Benchmarking}}
\author{Albert Mathews \\ Independent Researcher}
\date{\today}
\begin{document}
\maketitle
\begin{abstract}
my abstract here
\end{abstract}
\section{Introduction}
The study benchmarks the performance of four AI models---ChatGPT, Grok, Gemini, and Claude---in detecting cytopathic effects (CPE) in Vero cell culture images. Each model is implemented in individual Python scripts (\texttt{individual_image_chatgpt.py}, \texttt{individual_image_grok.py}, \texttt{individual_image_gemini.py}, \texttt{individual_image_claude.py}) that process microscope images to identify CPE presence and types. These scripts use multimodal AI capabilities to analyze image features such as cell rounding, vacuolation, and detachment, comparing results against ground truth from a Contract Research Organization (CRO).
\section{Methods}
The methodology involves the following steps:
\begin{itemize}
\item Convert raw dataset images to PNGs of the same size for uniform input to AI models.
\item Run inference on each AI model using the respective Python scripts to generate JSON outputs with CPE detection results.
\item Compare the AI results against CRO ground truth to compute accuracies for CPE detection and type identification.
\end{itemize}
\end{document}